\section{Related Work}

\paragraph{Online platforms, networks, and privacy externalities.}
Online-platform research shows how shared infrastructure and cross-side interactions create network effects among heterogeneous participants \cite{rochet2003platform,armstrong2006competition,katz1985network}. Information-spread models similarly show how local transmissions can have system-level consequences in networked environments \cite{kempe2003maximizing}. Privacy economics studies disclosure incentives and the costs created when information is reused outside its original context \cite{acquisti2016economics}. FlowFence contributes an enforceable runtime mediation primitive for a related AI-platform setting: one principal's write to shared agent state can impose privacy exposure costs on downstream principals.

\paragraph{Strategic ML and LLM-agent security.}
Strategic-classification and performative-prediction work studies how participants adapt to learned systems and platform decisions \cite{hardt2016strategic,perdomo2020performative}. In LLM-agent systems, indirect prompt injection shows that untrusted data can steer tool-integrated applications \cite{greshake2023indirect}, and recent benchmarks formalize prompt-injection, tool-agent, and memory-poisoning risks \cite{liu2023promptinjection,debenedetti2024agentdojo,chen2024agentpoison,zhang2025asb}. FlowFence is complementary: it measures and mediates whether sensitive or poisoned content propagates through memories, workspaces, messages, tools, and external outputs.

\paragraph{Information-flow control and runtime enforcement.}
Classical protection principles and information-flow models provide the foundation for controlling how sensitive data moves through systems \cite{saltzer1975protection,denning1976lattice}. Later work develops language-based information-flow security, decentralized labels, declassification, dynamic taint analysis, and runtime enforcement \cite{myers1999jflow,sabelfeld2003language,zdancewic2002robust,newsome2005dynamic,schneider2000enforceable}. FlowFence adapts this intuition to LLM-agent runtimes, where payloads are natural-language artifacts, agents act for different principals, and useful disclosure often requires policy-safe abstraction rather than binary allow/deny labels.

\paragraph{Topology-aware and memory-centric multi-agent systems.}
Multi-agent LLM systems rely on shared memory, summaries, and collaboration topologies, which also shape safety risk. Agent Smith illustrates how local compromise can spread across interacting agents \cite{gu2024agentsmith}; G-Safeguard studies topology-aware guardrails for LLM-based multi-agent systems \cite{wang2025gsafeguard}; A-MEM studies agent memory mechanisms \cite{xu2025amem}; and recent work identifies privacy risks in agent memory and autonomous web agents \cite{wang2025mextra,zharmagambetov2025agentdam}. FlowFence focuses on platform-level privacy propagation: when raw content should be blocked, quarantined, or rewritten before it enters shared state.
